\documentclass[paper=a4, fontsize=10pt]{jlreq}

\usepackage{amsmath, amssymb}
\usepackage{enumerate}
\usepackage{tikz}
\usepackage{listings, xcolor}

\lstset{
  basicstyle = {\ttfamily},
  frame = {tbrl},
  breaklines = true,
  numbers = left,
  showspaces = false,
  showstringspaces = false,
  showtabs = false,
  keywordstyle = \color{blue},
  commentstyle = {\color[HTML]{1AB91A}},
  identifierstyle = \color{black},
  stringstyle = \color{brown},
  captionpos = t
}

\title{サイバーメディアとXR \\ Topic A: クラウドセキュリティとゼロトラストセキュリティ}
\author{学籍番号 1280391 \\ 細川 夏風}
\date{\today}

\begin{document}

\maketitle

\section{現状と課題}
クラウドセキュリティの現状は単純である．危機に瀕しているのである．理由は至って単純である．クラウドは相互接続された計算資源を分割して顧客に貸し出しているに過ぎない．それは見かけ上は接続のないデバイス同士でも物理的にはそれぞれは接続されている．それだけではない．今や大規模なシステム(クラウドシステムも含む)はシステム自体を複雑に接続し開発を行っている．しかし，なんの対処も行われていないわけではない．それぞれをシステムで分割し，できるだけ独立に近い形を保ちデバイス同士，システム同士の領域を侵さないようにしている．しかし，それにも限界があるのが現状である．クラウドシステムに限らずシステムの運用・保守，どの工程にも人間が介在している．それは，システムの設計には人間が必要不可欠だからである．それ故にセキュリティホールがある．クラウドシステムもその大部分をそれを構築しているシステムに依存しているが，それ自体の安全性は\texttt{KADOKAWA}の件から懐疑的だ．だが，その現状に(\texttt{AI}: 人工知能)という鋭い刃が入った．かに思われた．人間が完璧なシステムを理解していないために完璧なシステムが完成することはない．このような現状に加え，星の数ほど存在しているクラッカーとその対処に現状で行うことのできる暫定解として出現したのがゼロトラストセキュリティである．

\section{ゼロトラストセキュリティ}
ゼロトラストセキュリティとはその名のとおりである．誰も信じないという開発思想の元設計されたセキュリティモデルである．これまでは信じるものを決め，それ以外を排除，検証という形で通信，接続を行っていた．しかし，信用しているのデバイスの権限が奪取されてしまってはまるで意味をなさない．その実例が\texttt{TOYOTA}のセキュリティインシデントである．この攻撃はサプライチェーン攻撃に当てはまるものであり，部品メーカー企業の脆弱性のある\texttt{VPN}から最終的には\texttt{TOYOTA}のネットワークまで侵入された．これはサプライチェーン先という内部ネットワークを信じ，検証を行わなかったために起こったセキュリティインシデントの最たる例であると言える．このようなセキュティインシデントを限りなく減らせるのがゼロトラストセキュリティである．理由は前述の通り，すべての要求に対して検証を行う．それが社内ネットワークでも自身の隣のデスクのデバイスからの要求であってもだ．検証はシステムで行われることがほとんどだ．もちろんこのシステムにも欠点は存在する．それは都度行われる検証がボトルネックになることと，これ自体を構築することにコストがかかることだ．どちらも納得に足る理由であると思われる．これまで必要のなかった部分にまで検証が行われるのだ．そこに割かれる計算コストはこれまでの$0$と比べるとより顕著に現れることは自明であるといえる．そして，このモデルの構築にかかるコストは，単純に現行のシステムから必要部分を洗い出し，それに対して適切な検証，適切な経路等を再度組み直すという大規模工事になる．それにかかる人的コストも構築コストもシステムの規模に応じて指数的になることは理解に易いと思われる．しかし，それでも導入するメリットが存在する．至って単純にそれはありとあらゆる情報を守ることができることである．現状，情報漏洩やクラッカーからの攻撃は社会的信用を失墜させる．信用が失墜したシステムが用いられるというケースはまず存在しない．すなわち，事業の死を意味することになる．そのため，コストがかかってもこれを導入する意義は大いに存在する．最近では市町村のシステムにこれが導入されるというケースもある．やはり公的なシステムは信用に割く割合は大きいことがわかる．

\section{考察}
ゼロトラストセキュリティが以下にシステムとして強力であるか，現行のシステムが以下に現代のセキュリティ事情に適していないかは前述の通り，理解の助けになったと思われる．私自身はもはやこのシステムモデル以外にクラウド上でも現行のシステム上でもセキュリティインシデントを減らすことはできないと考えている．セキュリティの天敵はほかでもなく人間である．\texttt{AI}の助けがあるとはいえ，システムのセキュリティホールを発見しその部分を攻撃することは非常に難しい．しかし，人間はとても騙しやすい．関係者を装ったメールを大量に送信すれば，何かのシステムは舵手で来てしまう可能性が高い．事実前述の\texttt{TOYOTA}の一件はフィッシングメールが引き金ではないかと言われている．最大のセキュリティホールである人間に対して，ある種それを信用しないという方法を取っているゼロトラストセキュリティがセキュリティの暫定解であると言えると考える．

\end{document}
